%% kapa: exam=1 round=35 year=2024 period=1 form=A subject=civil_law subject_ko=민법 seq=40 number=40
%% source: Q-Net 감정평가사 기출문제 (cst00309) artlSeq=5213229 file=raw/1008/5213229/1교시 문제지 원본.pdf
%% conversion: kordoc 4.24.51 -> tools/gen1.py (auto); review=proofread-sonnet; flags=none
\begin{question}{40}
甲은 乙에 대한 3억 원의 채권을 담보하기 위하여 乙소유 X토지와 丙소유 Y토지에 대하여 각각 1번 공동저당권을 취득하였고, 丁은 X에 대하여 피담보채권액 2억 원의 2번 저당권을 취득하였다. 그 후, 甲이 Y에 대한 경매를 신청하여 매각대금 2억 원을 배당받은 후 X에 대한 경매를 신청하여 X가 3억 원에 매각된 경우, 丁이 X의 매각대금에서 배당받을 수 있는 금액은? (경매비용ㆍ이자 등은 고려하지 않으며, 다툼이 있으면 판례에 따름)
\begin{choices}
  \item 0원
  \item 5천만 원
  \item 1억 원
  \item 1억 5천만 원
  \item 2억 원
\end{choices}
\end{question}
\answer{1}
